\documentclass[12pt,a4paper]{report}

% ===================================================================
% PACKAGES
% ===================================================================
\usepackage[utf8]{inputenc}
\usepackage[T1]{fontenc}
\usepackage[french]{babel}
\usepackage{graphicx}
\usepackage{float}
\usepackage{geometry}
\usepackage{hyperref}
\usepackage{caption}
\usepackage{booktabs}
\usepackage{listings}
\usepackage{xcolor}
\usepackage{titlesec}
\usepackage{fancyhdr}
\usepackage{longtable}
\usepackage{tikz}
\usepackage{array}
\usepackage{multirow}
\usepackage{enumitem}
\usepackage{tcolorbox}
\usepackage{tabularx}
\usepackage{colortbl}
\usepackage{lmodern}
\usepackage{microtype}
\usetikzlibrary{positioning, shapes, shapes.geometric, arrows.meta,
                fit, calc, shadows, backgrounds, decorations.pathreplacing}

% ===================================================================
% MISE EN PAGE
% ===================================================================
\geometry{top=2.5cm, bottom=2.5cm, left=3cm, right=2.5cm}

% ===================================================================
% COULEURS
% ===================================================================
\definecolor{medsysblue}{RGB}{37,99,235}
\definecolor{medsysdark}{RGB}{30,58,138}
\definecolor{medsysgreen}{RGB}{5,150,105}
\definecolor{medsysgray}{RGB}{107,114,128}
\definecolor{medsyslight}{RGB}{239,246,255}
\definecolor{entitybg}{RGB}{219,234,254}
\definecolor{relbg}{RGB}{253,230,138}
\definecolor{attrbg}{RGB}{248,250,252}
\definecolor{codebg}{RGB}{248,249,250}
\definecolor{codegreen}{RGB}{22,101,52}
\definecolor{codeblue}{RGB}{29,78,216}
\definecolor{codered}{RGB}{185,28,28}
\definecolor{colegray}{RGB}{107,114,128}
\definecolor{tabhead}{RGB}{30,58,138}
\definecolor{roweven}{RGB}{239,246,255}

% ===================================================================
% LISTINGS
% ===================================================================
\lstdefinestyle{sql}{
  language=SQL,
  backgroundcolor=\color{codebg},
  commentstyle=\color{colegray}\itshape,
  keywordstyle=\color{codeblue}\bfseries,
  stringstyle=\color{codered},
  basicstyle=\ttfamily\small,
  breaklines=true,
  frame=single, framesep=4pt, rulecolor=\color{gray!30},
  numbers=left, numberstyle=\tiny\color{colegray}, numbersep=6pt,
  tabsize=2, showstringspaces=false,
  morekeywords={AUTO\_INCREMENT,ENUM,UNIQUE,DEFAULT,REFERENCES,
                ON,DELETE,CASCADE,ENGINE,CHARSET,COLLATE,COMMENT}
}
\lstdefinestyle{java}{
  language=Java,
  backgroundcolor=\color{codebg},
  commentstyle=\color{colegray}\itshape,
  keywordstyle=\color{codeblue}\bfseries,
  stringstyle=\color{codered},
  basicstyle=\ttfamily\small,
  breaklines=true,
  frame=single, framesep=4pt, rulecolor=\color{gray!30},
  numbers=left, numberstyle=\tiny\color{colegray}, numbersep=6pt,
  tabsize=2
}
\lstdefinestyle{plantuml}{
  backgroundcolor=\color{codebg},
  basicstyle=\ttfamily\footnotesize,
  breaklines=true,
  frame=single, framesep=4pt, rulecolor=\color{gray!30},
  tabsize=2, showstringspaces=false
}
\lstset{style=sql}

% ===================================================================
% EN-TETE / PIED DE PAGE
% ===================================================================
\pagestyle{fancy}
\fancyhf{}
\fancyhead[L]{\small\textcolor{medsysgray}{\leftmark}}
\fancyhead[R]{\small\textbf{\textcolor{medsysblue}{MedSys}}}
\fancyfoot[C]{\textcolor{medsysgray}{\thepage}}
\renewcommand{\headrulewidth}{0.4pt}
\renewcommand{\headrule}{\hbox to\headwidth{\color{medsysblue!30}\leaders\hrule height \headrulewidth\hfill}}

% ===================================================================
% TITRES
% ===================================================================
\titleformat{\chapter}[block]
  {\normalfont\Huge\bfseries\color{medsysdark}}
  {\thechapter.}{0.8em}{}[\vspace{2pt}\color{medsysblue}\hrule]
\titlespacing*{\chapter}{0pt}{-10pt}{20pt}

\titleformat{\section}
  {\normalfont\Large\bfseries\color{medsysblue}}
  {\thesection}{0.8em}{}
\titleformat{\subsection}
  {\normalfont\large\bfseries\color{medsysgray!80!black}}
  {\thesubsection}{0.8em}{}
\titleformat{\subsubsection}
  {\normalfont\normalsize\bfseries\color{medsysgray}}
  {\thesubsubsection}{0.8em}{}

% ===================================================================
% HYPERLINKS
% ===================================================================
\hypersetup{
  colorlinks=true, linkcolor=medsysblue,
  urlcolor=medsysblue, citecolor=medsysblue,
  pdftitle={MedSys -- Rapport de Projet},
  pdfauthor={Équipe MedSys},
  pdfsubject={Système de Gestion Hospitalière -- Architecture Microservices}
}

% ===================================================================
% TIKZ STYLES GLOBAUX
% ===================================================================
\tikzset{
  % MCD -- entité
  mcde/.style={
    draw=medsysblue, thick, rectangle, rounded corners=3pt,
    fill=entitybg, minimum width=3.6cm, minimum height=0.75cm,
    align=center, font=\small\bfseries\color{medsysdark},
    drop shadow={shadow xshift=1pt, shadow yshift=-1pt, opacity=0.2}
  },
  % MCD -- relation
  mcdr/.style={
    draw=orange!70!red, thick, diamond, fill=relbg,
    minimum width=2.8cm, aspect=2.2,
    align=center, font=\scriptsize\bfseries\color{orange!60!black}
  },
  % MCD -- cardinalité
  card/.style={font=\scriptsize\bfseries, color=medsysblue!80!black},
  % Architecture -- service box
  svcbox/.style={
    draw, thick, rounded corners=5pt, fill=#1,
    minimum width=3.2cm, minimum height=1.2cm,
    align=center, font=\small\bfseries
  },
  % Flèche architecture
  svcline/.style={-Stealth, thick, color=medsysgray},
  svcdbline/.style={-Stealth, thick, dashed, color=medsysgray!60}
}

% ===================================================================
% TCOLORBOX
% ===================================================================
\tcbuselibrary{skins, breakable}
\newtcolorbox{infobox}[1][]{
  enhanced, breakable,
  colback=medsyslight, colframe=medsysblue,
  fonttitle=\bfseries\small, title={#1},
  leftrule=4pt, right=4pt, top=4pt, bottom=4pt,
  arc=3pt
}
\newtcolorbox{notebox}{
  enhanced, colback=yellow!10, colframe=orange!70,
  leftrule=4pt, arc=3pt, fonttitle=\bfseries\small
}

% ===================================================================
% COMMANDES UTILITAIRES
% ===================================================================
\newcommand{\port}[1]{\texttt{\small\textcolor{medsysblue}{#1}}}
\newcommand{\tech}[1]{\textbf{\textcolor{medsysdark}{#1}}}
\newcommand{\code}[1]{\texttt{\small #1}}

% Tableau centré bleu
\newcolumntype{H}{>{\columncolor{tabhead}\color{white}\bfseries\small}c}
\newcolumntype{E}{>{\columncolor{roweven}}l}

% ===================================================================
% DOCUMENT
% ===================================================================
\begin{document}

% -------------------------------------------------------------------
% PAGE DE GARDE
% -------------------------------------------------------------------
\begin{titlepage}
  \begin{tikzpicture}[remember picture, overlay]
    % Fond bleu en haut
    \fill[medsysdark] (current page.north west) rectangle
      ([yshift=-9cm]current page.north east);
    % Bande bleue claire
    \fill[medsysblue] ([yshift=-9cm]current page.north west) rectangle
      ([yshift=-9.3cm]current page.north east);
    % Cercles décoratifs
    \fill[white, opacity=0.04]
      ([xshift=5cm, yshift=-4cm]current page.north west) circle (4cm);
    \fill[white, opacity=0.04]
      ([xshift=16cm, yshift=-2cm]current page.north west) circle (6cm);
  \end{tikzpicture}

  \vspace*{1.5cm}
  \centering\color{white}
  {\fontsize{64}{70}\selectfont\bfseries MedSys\par}
  \vspace{0.4cm}
  {\Large\bfseries Système de Gestion Hospitalière\par}
  \vspace{0.25cm}
  {\normalsize Architecture Microservices \textbullet{} Spring Boot \textbullet{} ASP.NET Core \textbullet{} React\par}
  \vspace{0.8cm}

  \color{black}
  \vspace{3.5cm}

  \begin{tcolorbox}[enhanced, colback=white, colframe=medsysblue!30,
    width=14cm, arc=8pt, boxrule=1pt,
    drop shadow={opacity=0.15}]
    \centering
    {\Large\bfseries\color{medsysdark} Rapport de Conception\par}
    \vspace{0.3cm}
    {\normalsize Modélisation Merise (MCD / MLD / MPD)\\[4pt]
     Diagrammes UML (Cas d'utilisation, Classes, Séquences, Composants)\\[4pt]
     Description technique complète\par}
  \end{tcolorbox}

  \vspace{1.5cm}
  \begin{tabular}{cc}
    \toprule
    \textbf{Module} & \textbf{Technologie} \\
    \midrule
    ms-auth & Java 17 · Spring Boot 3 · MySQL \\
    ms-patient-personnel & Java 17 · Spring Boot 3 · MySQL \\
    MedicalAppointments & ASP.NET Core 8 · SQL Server \\
    ms-personnel & ASP.NET Core 8 · MongoDB \\
    medsys-web & React 18 · Vite · Axios \\
    \bottomrule
  \end{tabular}

  \vfill
  {\small\color{medsysgray} Rapport généré le \today\par}
\end{titlepage}

% -------------------------------------------------------------------
% RÉSUMÉ
% -------------------------------------------------------------------
\chapter*{Résumé}
\addcontentsline{toc}{chapter}{Résumé}

\begin{infobox}[MedSys en bref]
\textbf{MedSys} est une plateforme hospitalière complète basée sur une architecture
microservices hétérogène, intégrant des technologies Java (Spring Boot), .NET
(ASP.NET Core) et un frontend React unifié. Elle couvre l'ensemble du parcours
patient : authentification sécurisée, dossier médical électronique, prise de
rendez-vous, gestion du personnel et notifications en temps réel.
\end{infobox}

\vspace{0.5cm}

\noindent Le système est organisé autour de cinq modules principaux :

\begin{itemize}[leftmargin=*, itemsep=4pt]
  \item \tech{ms-auth} (port 8082) : Authentification JWT avec refresh token, protection
    brute-force, RBAC à 5 rôles (PATIENT, MÉDECIN, PERSONNEL, ADMIN, DIRECTEUR),
    audit logs, vérification email.

  \item \tech{ms-patient-personnel} (port 8081) : Dossier médical électronique complet —
    consultations, ordonnances, analyses de laboratoire, radiologies, hospitalisations,
    certificats médicaux, messagerie médecin-patient, export PDF, QR Code, tableau de bord direction.

  \item \tech{MedicalAppointments} (port 5000) : Prise de rendez-vous en ligne (patient
    anonyme ou connecté), gestion des créneaux par le médecin (vue semaine, création en masse,
    blocage), liste d'attente automatique avec notification email, SignalR temps réel.

  \item \tech{ms-personnel} (port variable) : Gestion des ressources humaines hospitalières —
    médecins, infirmiers, personnels, plannings, absences, équipes, disponibilités,
    demandes de modification. Base de données MongoDB.

  \item \tech{medsys-web} (port 5173) : Interface React unifiée connectée à tous les
    microservices via Vite proxy. Tableaux de bord par rôle, prise de RDV intégrée,
    statistiques direction.
\end{itemize}

\vspace{0.5cm}
\noindent\textbf{Mots-clés :} Microservices, Spring Boot, ASP.NET Core, React, JWT, MySQL, MongoDB,
RabbitMQ, SignalR, Merise, UML.

% -------------------------------------------------------------------
% TABLE DES MATIÈRES
% -------------------------------------------------------------------
\tableofcontents
\newpage

% -------------------------------------------------------------------
% CHAPITRES
% -------------------------------------------------------------------
\chapter{Introduction et Contexte}

\section{Contexte du projet}

Le secteur hospitalier fait face à une transformation numérique profonde. La gestion
manuelle des dossiers patients, la prise de rendez-vous par téléphone et la coordination
entre les différents services génèrent des inefficacités coûteuses et des risques d'erreurs
médicales. \textbf{MedSys} répond à ce besoin en proposant une plateforme numérique
intégrée, pensée pour les établissements de santé de taille moyenne à grande.

\section{Objectifs}

\begin{itemize}[itemsep=4pt]
  \item \textbf{Centraliser} les dossiers médicaux électroniques dans un système sécurisé
    et accessible aux professionnels de santé autorisés.
  \item \textbf{Simplifier} la prise de rendez-vous pour les patients, avec ou sans compte,
    grâce à une interface web intuitive.
  \item \textbf{Automatiser} les notifications (confirmations, rappels J-1, libération
    de créneaux) par email et en temps réel.
  \item \textbf{Outiller} la direction avec des tableaux de bord analytiques complets.
  \item \textbf{Gérer} les ressources humaines du personnel hospitalier (plannings,
    absences, équipes).
  \item \textbf{Sécuriser} l'ensemble des accès via JWT, RBAC et audit logs.
\end{itemize}

\section{Périmètre fonctionnel}

\begin{table}[H]
\centering
\renewcommand{\arraystretch}{1.4}
\begin{tabularx}{\textwidth}{|l|X|l|}
\hline
\rowcolor{tabhead}\color{white}\textbf{Module} &
\color{white}\textbf{Fonctionnalités principales} &
\color{white}\textbf{Utilisateurs} \\
\hline
ms-auth & Inscription, connexion, JWT, refresh token, RBAC, reset mdp, audit &
Patient, Admin \\
\hline
\rowcolor{roweven}
ms-patient-personnel & Dossier médical, consultations, ordonnances, analyses, PDF,
QR Code, messagerie, stats & Patient, Médecin, Direction \\
\hline
MedicalAppointments & Créneaux, réservation, liste d'attente, notifications, SignalR &
Patient, Médecin, Secrétaire \\
\hline
\rowcolor{roweven}
ms-personnel & Plannings, absences, équipes, disponibilités, demandes &
Personnel, Direction \\
\hline
medsys-web & Interface unifiée, tableaux de bord par rôle, prise de RDV &
Tous les rôles \\
\hline
\end{tabularx}
\caption{Périmètre fonctionnel de MedSys}
\end{table}

\section{Technologies utilisées}

\begin{table}[H]
\centering
\renewcommand{\arraystretch}{1.4}
\begin{tabular}{|l|l|l|}
\hline
\rowcolor{tabhead}\color{white}\textbf{Couche} &
\color{white}\textbf{Technologie} &
\color{white}\textbf{Version} \\
\hline
Backend Java & Spring Boot, Spring Security, Spring Data JPA & 3.x \\
\rowcolor{roweven}
Backend .NET & ASP.NET Core, Entity Framework Core & 8.0 \\
Frontend web & React, Vite, Axios & 18 / 5.x \\
\rowcolor{roweven}
Frontend mobile & React Native (Expo) & SDK 50 \\
Base de données & MySQL, SQL Server, MongoDB & 8 / 2022 / 7 \\
\rowcolor{roweven}
Messagerie async & RabbitMQ & 3.x \\
Temps réel & SignalR (.NET) & 8.0 \\
\rowcolor{roweven}
Auth & JWT (HS256), Refresh Token & -- \\
Conteneurisation & Docker (Dockerfile par service) & 24.x \\
\hline
\end{tabular}
\caption{Stack technologique de MedSys}
\end{table}

\section{Méthodes de modélisation}

Ce rapport présente la modélisation du système selon deux approches complémentaires :

\begin{itemize}[itemsep=3pt]
  \item \textbf{Merise} : modélisation des données à travers trois niveaux —
    MCD (conceptuel), MLD (logique) et MPD (physique). Adapté à la présentation
    de la structure des bases de données relationnelles.
  \item \textbf{UML 2.x} : modélisation comportementale et structurelle — diagrammes
    de cas d'utilisation, de classes, de séquences et de composants.
\end{itemize}

\chapter{Architecture du Système}

\section{Vue d'ensemble}

MedSys adopte une architecture \textbf{microservices hétérogène} : les services
communiquent via REST (synchrone) et RabbitMQ (asynchrone). Chaque service possède
sa propre base de données (principe \textit{database per service}).

\begin{figure}[H]
\centering
\begin{tikzpicture}[node distance=1.4cm and 1.6cm]

  % ---- Couche Frontend ----
  \node[draw, thick, rounded corners=6pt, fill=medsysgreen!15,
        minimum width=11cm, minimum height=1.1cm,
        align=center, font=\small\bfseries] (web)
    {medsys-web \quad \textcolor{medsysgray}{\small(React + Vite, port 5173)}};

  % ---- Microservices ----
  \node[draw, thick, rounded corners=5pt, fill=entitybg,
        minimum width=2.5cm, minimum height=1.1cm, align=center,
        font=\small\bfseries, below left=1.6cm and 3.2cm of web] (auth)
    {ms-auth\\[-2pt]\textcolor{medsysgray}{\tiny port 8082}};

  \node[draw, thick, rounded corners=5pt, fill=entitybg,
        minimum width=2.5cm, minimum height=1.1cm, align=center,
        font=\small\bfseries, below=1.6cm of web] (patient)
    {ms-patient\\[-2pt]\textcolor{medsysgray}{\tiny port 8081}};

  \node[draw, thick, rounded corners=5pt, fill=relbg!60,
        minimum width=2.5cm, minimum height=1.1cm, align=center,
        font=\small\bfseries, below right=1.6cm and 0.4cm of web] (rdv)
    {Medical\\[-2pt]Appointments\\[-2pt]\textcolor{medsysgray}{\tiny port 5000}};

  \node[draw, thick, rounded corners=5pt, fill=medsysgreen!20,
        minimum width=2.5cm, minimum height=1.1cm, align=center,
        font=\small\bfseries, right=0.8cm of rdv] (personnel)
    {ms-personnel\\[-2pt]\textcolor{medsysgray}{\tiny port variable}};

  % ---- Bases de données ----
  \node[draw, thick, cylinder, shape border rotate=90,
        fill=gray!15, minimum width=1.6cm, minimum height=0.9cm,
        align=center, font=\tiny\bfseries,
        below=1.2cm of auth] (dbauth)
    {MySQL\\ms\_auth\_db};

  \node[draw, thick, cylinder, shape border rotate=90,
        fill=gray!15, minimum width=1.6cm, minimum height=0.9cm,
        align=center, font=\tiny\bfseries,
        below=1.2cm of patient] (dbpatient)
    {MySQL\\ms\_patient\_db};

  \node[draw, thick, cylinder, shape border rotate=90,
        fill=gray!15, minimum width=1.8cm, minimum height=0.9cm,
        align=center, font=\tiny\bfseries,
        below=1.2cm of rdv] (dbrdv)
    {SQL Server\\medical\_appt};

  \node[draw, thick, cylinder, shape border rotate=90,
        fill=gray!15, minimum width=1.6cm, minimum height=0.9cm,
        align=center, font=\tiny\bfseries,
        below=1.2cm of personnel] (dbpers)
    {MongoDB\\personnel\_db};

  % ---- RabbitMQ ----
  \node[draw, thick, rounded corners=4pt, fill=orange!20,
        minimum width=2.2cm, minimum height=0.8cm,
        align=center, font=\small\bfseries,
        below=1.2cm of dbauth, xshift=1.5cm] (rmq)
    {RabbitMQ\\[-2pt]\textcolor{medsysgray}{\tiny port 5672}};

  % ---- Flèches Frontend -> Services ----
  \draw[-Stealth, thick, medsysblue] (web.south) -- ++(0,-0.5) -|
    node[near start, above, font=\tiny] {/api} (auth.north);
  \draw[-Stealth, thick, medsysblue] (web.south) -- ++(0,-0.5) -|
    node[near start, above, font=\tiny] {/api/v1} (patient.north);
  \draw[-Stealth, thick, orange!70!red] (web.south) -- ++(0,-0.5) -|
    node[near start, above, font=\tiny] {/api/appts} (rdv.north);
  \draw[-Stealth, thick, medsysgreen!60!black] (web.south) -- ++(0,-0.5) -|
    (personnel.north);

  % ---- Services -> BDD ----
  \draw[-Stealth, thick, gray] (auth.south) -- (dbauth.north);
  \draw[-Stealth, thick, gray] (patient.south) -- (dbpatient.north);
  \draw[-Stealth, thick, gray] (rdv.south) -- (dbrdv.north);
  \draw[-Stealth, thick, gray] (personnel.south) -- (dbpers.north);

  % ---- RabbitMQ ----
  \draw[-Stealth, thick, dashed, orange!70] (auth.south) |-
    node[near end, right, font=\tiny] {events} (rmq.west);
  \draw[-Stealth, thick, dashed, orange!70] (patient.south) |-
    (rmq.north);

  % ---- Inter-service ----
  \draw[Stealth-Stealth, thick, dashed, medsysgray!70] (auth.east) --
    node[above, font=\tiny] {JWT verify} (patient.west);

\end{tikzpicture}
\caption{Architecture globale de MedSys}
\label{fig:architecture}
\end{figure}

\section{Flux de communication}

\subsection{Communication synchrone (REST/HTTP)}

Tous les échanges frontend $\rightarrow$ backend et inter-services urgents
utilisent REST sur HTTP/1.1. Le frontend Vite proxifie les requêtes :

\begin{table}[H]
\centering
\renewcommand{\arraystretch}{1.3}
\begin{tabular}{|l|l|l|}
\hline
\rowcolor{tabhead}\color{white}\textbf{Préfixe URL} &
\color{white}\textbf{Service cible} &
\color{white}\textbf{Port} \\
\hline
\code{/api/v1/auth/**} & ms-auth & 8082 \\
\rowcolor{roweven}
\code{/api/v1/patient/**} & ms-patient-personnel & 8081 \\
\code{/api/v1/patients/**} & ms-patient-personnel & 8081 \\
\rowcolor{roweven}
\code{/api/v1/directeur/**} & ms-patient-personnel & 8081 \\
\code{/api/appointments/**} & MedicalAppointments & 5000 \\
\rowcolor{roweven}
\code{/api/slots/**} & MedicalAppointments & 5000 \\
\code{/api/services/**} & MedicalAppointments & 5000 \\
\rowcolor{roweven}
\code{/api/waiting-list/**} & MedicalAppointments & 5000 \\
\code{/appointmenthub} & MedicalAppointments (SignalR) & 5000 \\
\hline
\end{tabular}
\caption{Table de routage Vite proxy}
\end{table}

\subsection{Communication asynchrone (RabbitMQ)}

Les événements métier critiques (inscription patient, création de RDV,
annulation) sont publiés sur l'exchange \code{medsys.exchange} de RabbitMQ.
Les consommateurs s'abonnent aux queues correspondantes :

\begin{table}[H]
\centering
\renewcommand{\arraystretch}{1.3}
\begin{tabular}{|l|l|l|}
\hline
\rowcolor{tabhead}\color{white}\textbf{Événement} &
\color{white}\textbf{Producteur} &
\color{white}\textbf{Consommateur} \\
\hline
\code{patient.registered} & ms-auth & ms-patient-personnel \\
\rowcolor{roweven}
\code{appointment.created} & ms-patient-personnel & -- \\
\code{appointment.cancelled} & ms-patient-personnel & -- \\
\rowcolor{roweven}
\code{notification.send} & ms-patient-personnel & -- \\
\hline
\end{tabular}
\caption{Événements RabbitMQ}
\end{table}

\subsection{Temps réel (SignalR)}

Le service MedicalAppointments expose un hub SignalR sur
\code{/appointmenthub}. Lorsqu'un créneau est libéré (annulation),
tous les clients connectés reçoivent une notification en temps réel,
déclenchant un rafraîchissement automatique du calendrier.

\section{Sécurité transversale}

\begin{itemize}[itemsep=4pt]
  \item \textbf{JWT HS256} : token d'accès (15 min) + refresh token (7 jours),
    stocké en \code{sessionStorage} côté client.
  \item \textbf{RBAC} : 5 rôles hiérarchiques — PATIENT, MÉDECIN, PERSONNEL,
    ADMIN, DIRECTEUR. Chaque endpoint est protégé par \code{@PreAuthorize}
    (Spring) ou \code{[Authorize(Roles="...")]} (.NET).
  \item \textbf{Brute-force protection} : \code{LoginRateLimiter} dans ms-auth,
    blocage après 5 tentatives échouées (fenêtre 15 min).
  \item \textbf{Audit logs} : chaque action sensible (connexion, création de
    compte, changement de mot de passe) est tracée dans la table
    \code{audit\_logs}.
  \item \textbf{CORS} : origines autorisées limitées à
    \code{localhost:5173} et \code{localhost:3000}.
\end{itemize}

\chapter{Modélisation Merise}

\section{Introduction}

La méthode Merise structure la modélisation des données en trois niveaux d'abstraction.
MedSys étant basé sur une architecture microservices, chaque service possède son propre
modèle de données indépendant.

\begin{table}[H]
\centering
\renewcommand{\arraystretch}{1.3}
\begin{tabular}{|l|l|l|}
\hline
\rowcolor{tabhead}\color{white}\textbf{Niveau} &
\color{white}\textbf{Modèle} &
\color{white}\textbf{Objectif} \\
\hline
Conceptuel & MCD & Quoi stocker ? Entités et associations métier \\
\rowcolor{roweven}
Logique    & MLD & Comment organiser ? Tables et clés étrangères \\
Physique   & MPD & Comment implémenter ? SQL CREATE TABLE \\
\hline
\end{tabular}
\caption{Les trois niveaux de la méthode Merise}
\end{table}

\section{Modèle Conceptuel de Données (MCD)}

\subsection{MCD -- ms-auth}

\begin{figure}[H]
\centering
\begin{tikzpicture}

  \node[mcde] (ua) {USER\_ACCOUNT};
  \node[draw, fill=attrbg, below=0 of ua, text width=3.6cm,
        align=left, font=\tiny, inner sep=5pt] {
    \underline{id} : INT (PK)\\
    email : VARCHAR(200)\\
    passwordHash : VARCHAR\\
    nom, prenom, cin\\
    role : ENUM(5 rôles)\\
    enabled, emailVerified\\
    refreshToken\\
    patientId, personnelId\\
    createdAt : DATETIME
  };

  \node[mcde, right=5.5cm of ua] (al) {AUDIT\_LOG};
  \node[draw, fill=attrbg, below=0 of al, text width=3.6cm,
        align=left, font=\tiny, inner sep=5pt] {
    \underline{id} : INT (PK)\\
    userId : INT (FK)\\
    action : VARCHAR(100)\\
    detail : TEXT\\
    ipAddress : VARCHAR(45)\\
    createdAt : DATETIME
  };

  \node[mcdr] (gen) at ([xshift=2.8cm]ua.east) {GÉNÈRE};

  \draw[thick, medsysblue] (ua.east) -- (gen.west)
    node[midway, above, card] {0,1};
  \draw[thick, medsysblue] (gen.east) -- (al.west)
    node[midway, above, card] {0,n};

\end{tikzpicture}
\caption{MCD -- ms-auth}
\end{figure}

\noindent\textbf{Description :} \texttt{USER\_ACCOUNT} centralise tous les comptes
quels que soient les rôles. Les \texttt{AUDIT\_LOG} tracent les actions sensibles
(connexion, reset mot de passe, création de compte).

\subsection{MCD -- ms-patient-personnel}

\begin{figure}[H]
\centering
\begin{tikzpicture}[every node/.style={font=\scriptsize}]

  \node[mcde] (pat) {PATIENT};
  \node[draw, fill=attrbg, below=0 of pat, text width=3cm,
        align=left, font=\tiny, inner sep=4pt] {
    \underline{id}, cin, nom, prenom\\
    dateNaissance, sexe\\
    groupeSanguin, telephone\\
    email, ville, mutuelle
  };

  \node[mcde, right=3.5cm of pat] (dos) {DOSSIER\_MEDICAL};
  \node[draw, fill=attrbg, below=0 of dos, text width=3cm,
        align=left, font=\tiny, inner sep=4pt] {
    \underline{id}, patientId (FK)\\
    numeroDossier\\
    dateCreation, observations
  };

  \node[mcde, right=3.5cm of dos] (con) {CONSULTATION};
  \node[draw, fill=attrbg, below=0 of con, text width=3cm,
        align=left, font=\tiny, inner sep=4pt] {
    \underline{id}, dossierId (FK)\\
    medecinId (FK)\\
    dateConsultation, motif\\
    diagnostic, traitement\\
    poids, temperature
  };

  \node[mcde, below=4cm of dos] (ord) {ORDONNANCE};
  \node[draw, fill=attrbg, below=0 of ord, text width=3cm,
        align=left, font=\tiny, inner sep=4pt] {
    \underline{id}, dossierId (FK)\\
    medecinId (FK)\\
    dateOrdonnance, type\\
    instructions
  };

  \node[mcde, below=4cm of pat] (doc) {DOCUMENT\_PATIENT};
  \node[draw, fill=attrbg, below=0 of doc, text width=3cm,
        align=left, font=\tiny, inner sep=4pt] {
    \underline{id}, patientId (FK)\\
    typeDocument\\
    nomFichier, taille\\
    dateUpload
  };

  \node[mcde, below=4cm of con] (ana) {ANALYSE\_LABO};
  \node[draw, fill=attrbg, below=0 of ana, text width=3cm,
        align=left, font=\tiny, inner sep=4pt] {
    \underline{id}, dossierId (FK)\\
    typeAnalyse, laboratoire\\
    statut, resultats\\
    dateAnalyse
  };

  \node[mcde, above=1.2cm of dos] (med) {MEDECIN};
  \node[draw, fill=attrbg, below=0 of med, text width=3cm,
        align=left, font=\tiny, inner sep=4pt] {
    \underline{id}, matricule\\
    nom, prenom, email\\
    specialiteId, serviceId
  };

  \node[mcdr] (poss) at ([xshift=1.75cm]pat.east) {POSSÈDE};
  \node[mcdr] (cont) at ([xshift=1.75cm]dos.east) {CONTIENT};
  \node[mcdr] (incl) at ([yshift=-2cm]dos.south) {INCLUT};
  \node[mcdr] (dep)  at ([yshift=-2cm]pat.south)  {DÉPOSE};
  \node[mcdr] (real) at ([yshift=0.5cm]dos.north) {RÉALISE};

  \draw[thick, medsysblue] (pat.east) -- (poss.west) node[midway,above,card]{1,1};
  \draw[thick, medsysblue] (poss.east) -- (dos.west) node[midway,above,card]{1,1};
  \draw[thick, medsysblue] (dos.east) -- (cont.west) node[midway,above,card]{1,1};
  \draw[thick, medsysblue] (cont.east) -- (con.west) node[midway,above,card]{0,n};
  \draw[thick, medsysblue] (dos.south) -- (incl.north) node[midway,left,card]{1,1};
  \draw[thick, medsysblue] (incl.east) -- (ana.west) node[midway,below,card]{0,n};
  \draw[thick, medsysblue] (incl.west) -- (ord.east) node[midway,below,card]{0,n};
  \draw[thick, medsysblue] (pat.south) -- (dep.north) node[midway,left,card]{1,1};
  \draw[thick, medsysblue] (dep.east) -- (doc.west) node[midway,below,card]{0,n};
  \draw[thick, medsysblue] (med.south) -- (real.north) node[midway,right,card]{0,n};
  \draw[thick, medsysblue] (real.south) -- (dos.north) node[midway,right,card]{0,n};

\end{tikzpicture}
\caption{MCD -- ms-patient-personnel (entités principales)}
\end{figure}

\begin{infobox}[Entités complémentaires non représentées]
Pour des raisons de lisibilité : \texttt{RADIOLOGIE}, \texttt{ANTECEDENT},
\texttt{HOSPITALISATION}, \texttt{CERTIFICAT\_MEDICAL}, \texttt{LIGNE\_ORDONNANCE},
\texttt{MESSAGE\_PATIENT}, \texttt{FAVORITE\_DOCTOR}, \texttt{SPECIALITE}, \texttt{SERVICE}.
\end{infobox}

\subsection{MCD -- MedicalAppointments}

\begin{figure}[H]
\centering
\begin{tikzpicture}[every node/.style={font=\scriptsize}]

  \node[mcde] (usr) {USERS};
  \node[draw, fill=attrbg, below=0 of usr, text width=3.2cm,
        align=left, font=\tiny, inner sep=4pt] {
    \underline{id}, fullName, email\\
    phone, passwordHash\\
    role : Doctor/Patient/Secretary\\
    isRegistered, penaltyUntil\\
    cancelCount, createdAt
  };

  \node[mcde, right=3.8cm of usr] (ts) {TIME\_SLOTS};
  \node[draw, fill=attrbg, below=0 of ts, text width=3.2cm,
        align=left, font=\tiny, inner sep=4pt] {
    \underline{id}, doctorId (FK)\\
    startTime, endTime\\
    status : Available/Reserved/\\
    \quad Cancelled/Blocked
  };

  \node[mcde, right=3.8cm of ts] (appt) {APPOINTMENTS};
  \node[draw, fill=attrbg, below=0 of appt, text width=3.2cm,
        align=left, font=\tiny, inner sep=4pt] {
    \underline{id}, timeSlotId (FK UNIQUE)\\
    patientId (FK), bookedById (FK)\\
    reason, status\\
    anonymousToken\\
    cancelledAt, cancelReason
  };

  \node[mcde, below=4cm of usr] (wl) {WAITING\_LIST};
  \node[draw, fill=attrbg, below=0 of wl, text width=3.2cm,
        align=left, font=\tiny, inner sep=4pt] {
    \underline{id}, doctorId (FK)\\
    weekStartDate : DATE\\
    patientName, email, phone\\
    notifiedAt, createdAt
  };

  \node[mcde, below=4cm of appt] (svc) {SERVICES};
  \node[draw, fill=attrbg, below=0 of svc, text width=3.2cm,
        align=left, font=\tiny, inner sep=4pt] {
    \underline{id}, name\\
    description, icon\\
    createdAt
  };

  \node[mcdr] (pos) at ([xshift=1.9cm]usr.east) {POSSÈDE};
  \node[mcdr] (conc) at ([xshift=1.9cm]ts.east) {CONCERNE};
  \node[mcdr] (ins) at ([yshift=-2cm]usr.south) {S'INSCRIT};
  \node[mcdr] (prop) at ([yshift=-2cm]ts.south) {PROPOSÉ PAR};

  \draw[thick, medsysblue] (usr.east) -- (pos.west) node[midway,above,card]{1,1};
  \draw[thick, medsysblue] (pos.east) -- (ts.west) node[midway,above,card]{0,n};
  \draw[thick, medsysblue] (ts.east) -- (conc.west) node[midway,above,card]{1,1};
  \draw[thick, medsysblue] (conc.east) -- (appt.west) node[midway,above,card]{0,1};
  \draw[thick, medsysblue] (usr.south) -- (ins.north) node[midway,left,card]{1,1};
  \draw[thick, medsysblue] (ins.east) -- (wl.west) node[midway,below,card]{0,n};
  \draw[thick, medsysblue] (usr.south east) -- (prop.north west) node[near start,right,card]{0,n};
  \draw[thick, medsysblue] (prop.south east) -- (svc.north west) node[near end,right,card]{0,n};

\end{tikzpicture}
\caption{MCD -- MedicalAppointments}
\end{figure}

\subsection{MCD -- ms-personnel}

\begin{figure}[H]
\centering
\begin{tikzpicture}[every node/.style={font=\scriptsize}]

  \node[mcde] (per) {PERSONNEL};
  \node[draw, fill=attrbg, below=0 of per, text width=3cm,
        align=left, font=\tiny, inner sep=4pt] {
    \underline{id} : ObjectId\\
    matricule, nom, prenom\\
    cin, email, telephone\\
    type : Médecin/Infirmier/...\\
    serviceAffecte
  };

  \node[mcde, right=3.8cm of per] (plan) {PLANNING};
  \node[draw, fill=attrbg, below=0 of plan, text width=3cm,
        align=left, font=\tiny, inner sep=4pt] {
    \underline{id} : ObjectId\\
    personnelId (ref)\\
    dateDebut, dateFin\\
    periodicite, statut
  };

  \node[mcde, right=3.8cm of plan] (cren) {CRENEAU};
  \node[draw, fill=attrbg, below=0 of cren, text width=3cm,
        align=left, font=\tiny, inner sep=4pt] {
    \underline{id} : ObjectId\\
    medecinId (ref)\\
    jourSemaine\\
    heureDebut, heureFin\\
    typeConsultation
  };

  \node[mcde, below=4cm of per] (abs) {ABSENCE};
  \node[draw, fill=attrbg, below=0 of abs, text width=3cm,
        align=left, font=\tiny, inner sep=4pt] {
    \underline{id} : ObjectId\\
    personnelId (ref)\\
    dateDebut, dateFin\\
    typeAbsence, statut, motif
  };

  \node[mcde, below=4cm of plan] (equi) {EQUIPE};
  \node[draw, fill=attrbg, below=0 of equi, text width=3cm,
        align=left, font=\tiny, inner sep=4pt] {
    \underline{id} : ObjectId\\
    nom, description\\
    responsableId (ref)\\
    membres : [ObjectId]
  };

  \node[mcde, below=4cm of cren] (dem) {DEMANDE\_MODIF};
  \node[draw, fill=attrbg, below=0 of dem, text width=3cm,
        align=left, font=\tiny, inner sep=4pt] {
    \underline{id} : ObjectId\\
    demandeurId (ref)\\
    typeModification\\
    statut, dateDemande
  };

  \node[mcdr] (a) at ([xshift=1.9cm]per.east) {A};
  \node[mcdr] (b) at ([xshift=1.9cm]plan.east) {COMPOSE};
  \node[mcdr] (c) at ([yshift=-2cm]per.south) {DÉCLARE};
  \node[mcdr] (d) at ([yshift=-2cm]plan.south) {APPARTIENT};
  \node[mcdr] (e) at ([yshift=-2cm]cren.south) {FAIT};

  \draw[thick, medsysblue] (per.east) -- (a.west) node[midway,above,card]{1,1};
  \draw[thick, medsysblue] (a.east) -- (plan.west) node[midway,above,card]{0,n};
  \draw[thick, medsysblue] (plan.east) -- (b.west) node[midway,above,card]{1,1};
  \draw[thick, medsysblue] (b.east) -- (cren.west) node[midway,above,card]{0,n};
  \draw[thick, medsysblue] (per.south) -- (c.north) node[midway,left,card]{1,1};
  \draw[thick, medsysblue] (c.east) -- (abs.west) node[midway,below,card]{0,n};
  \draw[thick, medsysblue] (per.south east) -- (d.north west) node[near start,right,card]{0,n};
  \draw[thick, medsysblue] (d.south west) -- (equi.north east) node[near end,right,card]{0,n};
  \draw[thick, medsysblue] (cren.south) -- (e.north) node[midway,left,card]{1,1};
  \draw[thick, medsysblue] (e.south) -- (dem.north) node[midway,left,card]{0,n};

\end{tikzpicture}
\caption{MCD -- ms-personnel}
\end{figure}

\section{Modèle Logique de Données (MLD)}

\subsection{MLD -- ms-auth}
\begin{itemize}[itemsep=2pt]
  \item \textbf{USER\_ACCOUNT}(\underline{id}, email, passwordHash, nom, prenom, cin, role, enabled, emailVerified, refreshToken, refreshTokenExpiry, patientId, personnelId, createdAt)
  \item \textbf{AUDIT\_LOG}(\underline{id}, \textit{\#userId} $\to$ USER\_ACCOUNT, action, detail, ipAddress, userAgent, createdAt)
\end{itemize}

\subsection{MLD -- ms-patient-personnel}
\begin{itemize}[itemsep=2pt]
  \item \textbf{PATIENT}(\underline{id}, cin, nom, prenom, dateNaissance, sexe, groupeSanguin, telephone, email, ville, mutuelle, numeroDossier)
  \item \textbf{MEDECIN}(\underline{id}, matricule, nom, prenom, \textit{\#specialiteId}, \textit{\#serviceId}, email)
  \item \textbf{DOSSIER\_MEDICAL}(\underline{id}, \textit{\#patientId} $\to$ PATIENT, numeroDossier, dateCreation)
  \item \textbf{CONSULTATION}(\underline{id}, \textit{\#dossierId}, \textit{\#medecinId}, dateConsultation, motif, diagnostic, traitement, poids, temperature, tensionSystolique, tensionDiastolique)
  \item \textbf{ORDONNANCE}(\underline{id}, \textit{\#dossierId}, \textit{\#medecinId}, dateOrdonnance, typeOrdonnance, instructions)
  \item \textbf{LIGNE\_ORDONNANCE}(\underline{id}, \textit{\#ordonnanceId}, medicament, dosage, dureeJours, posologie)
  \item \textbf{ANALYSE\_LABO}(\underline{id}, \textit{\#dossierId}, typeAnalyse, laboratoire, statut, resultats, dateAnalyse)
  \item \textbf{RADIOLOGIE}(\underline{id}, \textit{\#dossierId}, typeExamen, description, conclusion, dateExamen)
  \item \textbf{ANTECEDENT}(\underline{id}, \textit{\#dossierId}, typeAntecedent, description, dateDiagnostic, severite)
  \item \textbf{DOCUMENT\_PATIENT}(\underline{id}, \textit{\#patientId}, typeDocument, nomFichier, cheminFichier, taille, dateUpload)
  \item \textbf{MESSAGE\_PATIENT}(\underline{id}, \textit{\#patientId}, \textit{\#medecinId}, contenu, expediteur, lu, dateEnvoi)
  \item \textbf{FAVORITE\_DOCTOR}(\underline{id}, \textit{\#patientId}, \textit{\#medecinId}, dateAjout)
  \item \textbf{SPECIALITE}(\underline{id}, nom, description)
  \item \textbf{SERVICE}(\underline{id}, nom, description, \textit{\#chefServiceId})
\end{itemize}

\subsection{MLD -- MedicalAppointments}
\begin{itemize}[itemsep=2pt]
  \item \textbf{USERS}(\underline{id}, fullName, email, phone, passwordHash, role, isRegistered, penaltyUntil, cancelCount, createdAt)
  \item \textbf{TIME\_SLOTS}(\underline{id}, \textit{\#doctorId} $\to$ USERS, startTime, endTime, status, UNIQUE(doctorId, startTime))
  \item \textbf{APPOINTMENTS}(\underline{id}, \textit{\#timeSlotId}(UNIQUE), \textit{\#patientId} $\to$ USERS, \textit{\#bookedById} $\to$ USERS, reason, status, anonymousToken, cancelledAt, cancelReason)
  \item \textbf{WAITING\_LIST}(\underline{id}, \textit{\#doctorId}, weekStartDate, patientName, email, phone, notifiedAt, UNIQUE(doctorId, weekStartDate, email))
  \item \textbf{SERVICES}(\underline{id}, name, description, icon, createdAt)
  \item \textbf{DOCTOR\_SERVICES}(\textit{\#doctorId}, \textit{\#serviceId}, PK composite)
  \item \textbf{NOTIFICATIONS}(\underline{id}, recipientEmail, subject, body, sentAt, status)
  \item \textbf{AUDIT\_LOGS}(\underline{id}, \textit{\#userId}, action, entityType, entityId, detail, createdAt)
\end{itemize}

\section{Modèle Physique de Données (MPD)}

\subsection{MPD -- ms-auth (MySQL)}

\begin{lstlisting}[style=sql, caption={Tables ms-auth}]
CREATE TABLE user_accounts (
  id               BIGINT AUTO_INCREMENT PRIMARY KEY,
  email            VARCHAR(200) NOT NULL UNIQUE,
  password_hash    VARCHAR(255) NOT NULL,
  nom              VARCHAR(100),
  prenom           VARCHAR(100),
  cin              VARCHAR(20),
  role             ENUM('PATIENT','MEDECIN','PERSONNEL','ADMIN','DIRECTEUR')
                   NOT NULL DEFAULT 'PATIENT',
  enabled          BOOLEAN NOT NULL DEFAULT FALSE,
  email_verified   BOOLEAN NOT NULL DEFAULT FALSE,
  refresh_token    VARCHAR(512),
  refresh_token_expiry DATETIME,
  patient_id       BIGINT,
  personnel_id     BIGINT,
  created_at       DATETIME NOT NULL DEFAULT CURRENT_TIMESTAMP
) ENGINE=InnoDB DEFAULT CHARSET=utf8mb4;

CREATE TABLE audit_logs (
  id          BIGINT AUTO_INCREMENT PRIMARY KEY,
  user_id     BIGINT REFERENCES user_accounts(id) ON DELETE SET NULL,
  action      VARCHAR(100) NOT NULL,
  detail      TEXT,
  ip_address  VARCHAR(45),
  created_at  DATETIME NOT NULL DEFAULT CURRENT_TIMESTAMP
) ENGINE=InnoDB DEFAULT CHARSET=utf8mb4;
\end{lstlisting}

\subsection{MPD -- MedicalAppointments (SQL Server)}

\begin{lstlisting}[style=sql, caption={Tables MedicalAppointments}]
CREATE TABLE Users (
  Id           INT IDENTITY PRIMARY KEY,
  FullName     NVARCHAR(150) NOT NULL,
  Email        NVARCHAR(200) NOT NULL UNIQUE,
  Phone        NVARCHAR(20),
  PasswordHash NVARCHAR(255),
  Role         NVARCHAR(20) NOT NULL DEFAULT 'Patient',
  IsRegistered BIT NOT NULL DEFAULT 0,
  PenaltyUntil DATETIME2 NULL,
  CancelCount  INT NOT NULL DEFAULT 0,
  CreatedAt    DATETIME2 NOT NULL DEFAULT GETDATE()
);
CREATE TABLE TimeSlots (
  Id        INT IDENTITY PRIMARY KEY,
  DoctorId  INT NOT NULL REFERENCES Users(Id) ON DELETE CASCADE,
  StartTime DATETIME2 NOT NULL,
  EndTime   DATETIME2 NOT NULL,
  Status    NVARCHAR(20) NOT NULL DEFAULT 'Available',
  CONSTRAINT UQ_Doctor_Slot UNIQUE (DoctorId, StartTime)
);
CREATE TABLE Appointments (
  Id             INT IDENTITY PRIMARY KEY,
  TimeSlotId     INT NOT NULL UNIQUE REFERENCES TimeSlots(Id),
  PatientId      INT NOT NULL REFERENCES Users(Id),
  BookedById     INT NOT NULL REFERENCES Users(Id),
  Reason         NVARCHAR(500),
  Status         NVARCHAR(40) NOT NULL DEFAULT 'Confirmed',
  AnonymousToken NVARCHAR(100) NULL,
  CancelledAt    DATETIME2 NULL,
  CancelReason   NVARCHAR(500) NULL,
  CreatedAt      DATETIME2 NOT NULL DEFAULT GETDATE()
);
CREATE TABLE WaitingList (
  Id            INT IDENTITY PRIMARY KEY,
  DoctorId      INT NOT NULL REFERENCES Users(Id) ON DELETE CASCADE,
  WeekStartDate DATE NOT NULL,
  PatientName   NVARCHAR(150) NOT NULL,
  Email         NVARCHAR(200) NOT NULL,
  Phone         NVARCHAR(20),
  NotifiedAt    DATETIME2 NULL,
  CreatedAt     DATETIME2 NOT NULL DEFAULT GETDATE(),
  CONSTRAINT UQ_WaitingList UNIQUE (DoctorId, WeekStartDate, Email)
);
CREATE TABLE Services (
  Id   INT IDENTITY PRIMARY KEY,
  Name NVARCHAR(150) NOT NULL UNIQUE,
  Icon NVARCHAR(10)
);
CREATE TABLE DoctorServices (
  DoctorId INT NOT NULL REFERENCES Users(Id) ON DELETE CASCADE,
  ServiceId INT NOT NULL REFERENCES Services(Id) ON DELETE CASCADE,
  PRIMARY KEY (DoctorId, ServiceId)
);
\end{lstlisting}

\subsection{MPD -- ms-patient-personnel (MySQL, extrait)}

\begin{lstlisting}[style=sql, caption={Tables principales ms-patient-personnel}]
CREATE TABLE patients (
  id             BIGINT AUTO_INCREMENT PRIMARY KEY,
  cin            VARCHAR(20) NOT NULL UNIQUE,
  nom            VARCHAR(100) NOT NULL,
  prenom         VARCHAR(100) NOT NULL,
  date_naissance DATE,
  sexe           ENUM('MASCULIN','FEMININ') NOT NULL,
  groupe_sanguin VARCHAR(10),
  telephone      VARCHAR(20),
  email          VARCHAR(200),
  ville          VARCHAR(100),
  mutuelle       VARCHAR(100),
  numero_dossier VARCHAR(20) UNIQUE
) ENGINE=InnoDB DEFAULT CHARSET=utf8mb4;

CREATE TABLE dossiers_medicaux (
  id             BIGINT AUTO_INCREMENT PRIMARY KEY,
  patient_id     BIGINT NOT NULL UNIQUE
                 REFERENCES patients(id) ON DELETE CASCADE,
  numero_dossier VARCHAR(20) NOT NULL UNIQUE,
  date_creation  DATETIME NOT NULL DEFAULT CURRENT_TIMESTAMP,
  observations   TEXT
) ENGINE=InnoDB DEFAULT CHARSET=utf8mb4;

CREATE TABLE consultations (
  id                  BIGINT AUTO_INCREMENT PRIMARY KEY,
  dossier_id          BIGINT NOT NULL
                      REFERENCES dossiers_medicaux(id) ON DELETE CASCADE,
  medecin_id          BIGINT REFERENCES medecins(id),
  date_consultation   DATETIME NOT NULL,
  motif               TEXT,
  diagnostic          TEXT,
  traitement          TEXT,
  poids               DECIMAL(5,2),
  temperature         DECIMAL(4,1),
  tension_systolique  INT,
  tension_diastolique INT
) ENGINE=InnoDB DEFAULT CHARSET=utf8mb4;

CREATE TABLE ordonnances (
  id               BIGINT AUTO_INCREMENT PRIMARY KEY,
  dossier_id       BIGINT NOT NULL
                   REFERENCES dossiers_medicaux(id) ON DELETE CASCADE,
  medecin_id       BIGINT REFERENCES medecins(id),
  date_ordonnance  DATE NOT NULL,
  type_ordonnance  ENUM('SIMPLE','RENOUVELABLE','LONGUE_DUREE') NOT NULL,
  instructions     TEXT
) ENGINE=InnoDB DEFAULT CHARSET=utf8mb4;

CREATE TABLE lignes_ordonnance (
  id            BIGINT AUTO_INCREMENT PRIMARY KEY,
  ordonnance_id BIGINT NOT NULL
                REFERENCES ordonnances(id) ON DELETE CASCADE,
  medicament    VARCHAR(200) NOT NULL,
  dosage        VARCHAR(100),
  duree_jours   INT,
  posologie     TEXT
) ENGINE=InnoDB DEFAULT CHARSET=utf8mb4;

CREATE TABLE documents_patients (
  id            BIGINT AUTO_INCREMENT PRIMARY KEY,
  patient_id    BIGINT NOT NULL REFERENCES patients(id) ON DELETE CASCADE,
  type_document ENUM('ORDONNANCE','ANALYSE','RADIOLOGIE',
                     'CERTIFICAT','AUTRE') NOT NULL,
  nom_fichier   VARCHAR(255) NOT NULL,
  chemin        VARCHAR(512) NOT NULL,
  taille        BIGINT,
  date_upload   DATETIME NOT NULL DEFAULT CURRENT_TIMESTAMP
) ENGINE=InnoDB DEFAULT CHARSET=utf8mb4;
\end{lstlisting}

\chapter{Modélisation UML}

\section{Introduction}

La modélisation UML 2.x couvre quatre types de diagrammes pour MedSys :
les cas d'utilisation (fonctionnel), les classes (structurel),
les séquences (comportemental) et les composants (architectural).

\section{Diagrammes de Cas d'Utilisation}

\subsection{ms-auth -- Cas d'utilisation}

\begin{figure}[H]
\centering
\begin{tikzpicture}[every node/.style={font=\small}]

  % Acteurs
  \node[draw, circle, fill=medsyslight, minimum size=0.6cm] (pat) at (0,5) {};
  \node[below=0pt of pat, font=\tiny] {Patient};
  \node[draw, circle, fill=medsyslight, minimum size=0.6cm] (adm) at (0,1) {};
  \node[below=0pt of adm, font=\tiny] {Admin};
  \node[draw, circle, fill=medsyslight, minimum size=0.6cm] (any) at (0,3) {};
  \node[below=0pt of any, font=\tiny] {Tout utilisateur};

  % Système
  \draw[thick, medsysblue] (2,0) rectangle (11,7);
  \node[above, font=\small\bfseries\color{medsysblue}] at (6.5,7) {ms-auth};

  % Use cases
  \node[draw, ellipse, fill=entitybg, text width=3.2cm, align=center,
        font=\scriptsize] (uc1) at (7,6.3) {S'inscrire (Patient)};
  \node[draw, ellipse, fill=entitybg, text width=3.2cm, align=center,
        font=\scriptsize] (uc2) at (7,5.2) {Se connecter (JWT)};
  \node[draw, ellipse, fill=entitybg, text width=3.2cm, align=center,
        font=\scriptsize] (uc3) at (7,4.1) {Rafraîchir le token};
  \node[draw, ellipse, fill=entitybg, text width=3.2cm, align=center,
        font=\scriptsize] (uc4) at (7,3.0) {Réinitialiser le mot de passe};
  \node[draw, ellipse, fill=entitybg, text width=3.2cm, align=center,
        font=\scriptsize] (uc5) at (7,1.9) {Changer le mot de passe};
  \node[draw, ellipse, fill=relbg, text width=3.2cm, align=center,
        font=\scriptsize] (uc6) at (4,3.0) {Créer compte personnel/médecin};
  \node[draw, ellipse, fill=relbg, text width=3.2cm, align=center,
        font=\scriptsize] (uc7) at (4,1.5) {Activer / désactiver un compte};

  % Liens
  \draw[thick] (pat.east) -- (uc1.west);
  \draw[thick] (any.east) -- (uc2.west);
  \draw[thick] (any.east) -- (uc3.west);
  \draw[thick] (any.east) -- (uc4.west);
  \draw[thick] (any.east) -- (uc5.west);
  \draw[thick] (adm.east) -- (uc6.west);
  \draw[thick] (adm.east) -- (uc7.west);

\end{tikzpicture}
\caption{Diagramme de cas d'utilisation -- ms-auth}
\end{figure}

\subsection{ms-patient-personnel -- Cas d'utilisation}

\begin{figure}[H]
\centering
\begin{tikzpicture}[every node/.style={font=\scriptsize}]

  % Acteurs
  \node[draw, circle, fill=medsyslight, minimum size=0.5cm] (pat) at (0,8) {};
  \node[below=0pt of pat, font=\tiny] {Patient};
  \node[draw, circle, fill=medsyslight, minimum size=0.5cm] (med) at (0,4) {};
  \node[below=0pt of med, font=\tiny] {Médecin};
  \node[draw, circle, fill=relbg, minimum size=0.5cm] (dir) at (0,1) {};
  \node[below=0pt of dir, font=\tiny] {Directeur};

  \draw[thick, medsysblue] (1.5,0) rectangle (12,10);
  \node[above, font=\small\bfseries\color{medsysblue}] at (6.75,10) {ms-patient-personnel};

  % Patient UCs
  \foreach \y/\txt in {
    9.3/{Consulter son dossier médical},
    8.4/{Uploader un document},
    7.5/{Envoyer un message au médecin},
    6.6/{Prendre un rendez-vous},
    5.7/{Exporter PDF / QR Code},
    4.8/{Modifier son profil}
  }{
    \node[draw, ellipse, fill=entitybg, text width=3.5cm, align=center]
      (p\y) at (7.5,\y) {\txt};
    \draw[thick] (pat.east) -- (p\y.west);
  }

  % Médecin UCs
  \foreach \y/\txt in {
    4.0/{Voir liste patients},
    3.1/{Consulter dossier patient},
    2.2/{Envoyer message patient}
  }{
    \node[draw, ellipse, fill=entitybg, text width=3.5cm, align=center]
      (m\y) at (7.5,\y) {\txt};
    \draw[thick] (med.east) -- (m\y.west);
  }

  % Directeur UCs
  \foreach \y/\txt in {
    1.3/{Tableau de bord -- stats},
    0.4/{Exporter PDF d'un patient}
  }{
    \node[draw, ellipse, fill=relbg, text width=3.5cm, align=center]
      (d\y) at (7.5,\y) {\txt};
    \draw[thick] (dir.east) -- (d\y.west);
  }

\end{tikzpicture}
\caption{Diagramme de cas d'utilisation -- ms-patient-personnel}
\end{figure}

\subsection{MedicalAppointments -- Cas d'utilisation}

\begin{figure}[H]
\centering
\begin{tikzpicture}[every node/.style={font=\scriptsize}]

  \node[draw, circle, fill=medsyslight, minimum size=0.5cm] (pat) at (0,7) {};
  \node[below=0pt of pat, font=\tiny] {Patient (anonyme/connecté)};
  \node[draw, circle, fill=medsyslight, minimum size=0.5cm] (doc) at (0,3) {};
  \node[below=0pt of doc, font=\tiny] {Médecin};
  \node[draw, circle, fill=relbg, minimum size=0.5cm] (sec) at (0,1) {};
  \node[below=0pt of sec, font=\tiny] {Secrétaire};

  \draw[thick, medsysblue] (1.5,0) rectangle (12,9);
  \node[above, font=\small\bfseries\color{medsysblue}] at (6.75,9) {MedicalAppointments};

  \foreach \y/\txt/\col in {
    8.3/{Voir créneaux disponibles}/entitybg,
    7.4/{Réserver un créneau}/entitybg,
    6.5/{Annuler un rendez-vous}/entitybg,
    5.6/{S'inscrire liste d'attente}/entitybg,
    4.7/{Consulter mes rendez-vous}/entitybg
  }{
    \node[draw, ellipse, fill=\col, text width=3.8cm, align=center]
      (u\y) at (7.5,\y) {\txt};
    \draw[thick] (pat.east) -- (u\y.west);
  }
  \foreach \y/\txt in {
    3.6/{Créer/bloquer des créneaux},
    2.7/{Voir liste d'attente},
    1.8/{Voir tous les rendez-vous}
  }{
    \node[draw, ellipse, fill=entitybg, text width=3.8cm, align=center]
      (d\y) at (7.5,\y) {\txt};
    \draw[thick] (doc.east) -- (d\y.west);
  }
  \node[draw, ellipse, fill=relbg, text width=3.8cm, align=center]
    (s1) at (7.5,0.8) {Modifier un rendez-vous};
  \draw[thick] (sec.east) -- (s1.west);

\end{tikzpicture}
\caption{Diagramme de cas d'utilisation -- MedicalAppointments}
\end{figure}

\section{Diagrammes de Classes}

\subsection{Classes -- ms-auth}

\begin{figure}[H]
\centering
\begin{tikzpicture}[every node/.style={font=\scriptsize}]

  % UserAccount
  \node[draw, thick, fill=entitybg, text width=5.5cm] (ua) at (0,0) {
    \begin{center}\textbf{UserAccount}\end{center}
    \hrule\vspace{2pt}
    \texttt{- id : Long}\\
    \texttt{- email : String}\\
    \texttt{- passwordHash : String}\\
    \texttt{- nom, prenom, cin : String}\\
    \texttt{- role : Role}\\
    \texttt{- enabled, emailVerified : boolean}\\
    \texttt{- refreshToken : String}\\
    \texttt{- patientId, personnelId : Long}\\
    \vspace{2pt}\hrule\vspace{2pt}
    \texttt{+ isAccountNonExpired() : boolean}\\
    \texttt{+ getAuthorities() : Collection}
  };

  % Role (enum)
  \node[draw, thick, fill=relbg, text width=3.5cm] (role) at (7,1) {
    \begin{center}\textbf{\small$\ll$enumeration$\gg$\\Role}\end{center}
    \hrule\vspace{2pt}
    \texttt{PATIENT}\\
    \texttt{MEDECIN}\\
    \texttt{PERSONNEL}\\
    \texttt{ADMIN}\\
    \texttt{DIRECTEUR}
  };

  % AuditLog
  \node[draw, thick, fill=entitybg, text width=5.5cm] (al) at (0,-5) {
    \begin{center}\textbf{AuditLog}\end{center}
    \hrule\vspace{2pt}
    \texttt{- id : Long}\\
    \texttt{- userId : Long}\\
    \texttt{- action : String}\\
    \texttt{- detail : String}\\
    \texttt{- ipAddress : String}\\
    \texttt{- createdAt : LocalDateTime}
  };

  % JwtService
  \node[draw, thick, fill=attrbg, text width=5.5cm] (jwt) at (7,-3) {
    \begin{center}\textbf{JwtService}\end{center}
    \hrule\vspace{2pt}
    \texttt{- secretKey : String}\\
    \texttt{- expiration : long}\\
    \vspace{2pt}\hrule\vspace{2pt}
    \texttt{+ generateToken(user) : String}\\
    \texttt{+ generateRefreshToken(user)}\\
    \texttt{+ extractUsername(token) : String}\\
    \texttt{+ isTokenValid(token, user) : boolean}
  };

  % AuthService
  \node[draw, thick, fill=attrbg, text width=5.5cm] (as) at (7,-7) {
    \begin{center}\textbf{AuthService}\end{center}
    \hrule\vspace{2pt}
    \texttt{+ login(req) : AuthResponse}\\
    \texttt{+ register(req) : AuthResponse}\\
    \texttt{+ refreshToken(req) : AuthResponse}\\
    \texttt{+ resetPassword(req)}\\
    \texttt{+ changePassword(req)}
  };

  \draw[-Stealth, thick, medsysblue] (ua.east) -- (role.west)
    node[midway, above, font=\tiny] {utilise};
  \draw[-Stealth, thick, medsysblue] (ua.south) -- (al.north)
    node[midway, right, font=\tiny] {1 $\to$ 0..*};
  \draw[-Stealth, thick, dashed, medsysgray] (jwt.south) -- (as.north);
  \draw[-Stealth, thick, dashed, medsysgray] (as.west) -- (ua.east);

\end{tikzpicture}
\caption{Diagramme de classes -- ms-auth}
\end{figure}

\subsection{Classes -- MedicalAppointments (extrait)}

\begin{figure}[H]
\centering
\begin{tikzpicture}[every node/.style={font=\scriptsize}]

  \node[draw, thick, fill=entitybg, text width=4.8cm] (u) at (0,3) {
    \begin{center}\textbf{User}\end{center}
    \hrule\vspace{2pt}
    \texttt{+ Id : int}\\
    \texttt{+ FullName : string}\\
    \texttt{+ Email : string}\\
    \texttt{+ Role : UserRole}\\
    \texttt{+ IsRegistered : bool}\\
    \texttt{+ PenaltyUntil : DateTime?}\\
    \texttt{+ CancelCount : int}
  };

  \node[draw, thick, fill=entitybg, text width=4.8cm] (ts) at (6.5,3) {
    \begin{center}\textbf{TimeSlot}\end{center}
    \hrule\vspace{2pt}
    \texttt{+ Id : int}\\
    \texttt{+ DoctorId : int}\\
    \texttt{+ Doctor : User}\\
    \texttt{+ StartTime : DateTime}\\
    \texttt{+ EndTime : DateTime}\\
    \texttt{+ Status : SlotStatus}\\
    \texttt{+ Appointment : Appointment?}
  };

  \node[draw, thick, fill=entitybg, text width=4.8cm] (ap) at (6.5,-1) {
    \begin{center}\textbf{Appointment}\end{center}
    \hrule\vspace{2pt}
    \texttt{+ Id : int}\\
    \texttt{+ TimeSlotId : int}\\
    \texttt{+ PatientId : int}\\
    \texttt{+ BookedById : int}\\
    \texttt{+ Reason : string?}\\
    \texttt{+ Status : AppointmentStatus}\\
    \texttt{+ AnonymousToken : string?}
  };

  \node[draw, thick, fill=entitybg, text width=4.8cm] (wl) at (0,-1) {
    \begin{center}\textbf{WaitingListEntry}\end{center}
    \hrule\vspace{2pt}
    \texttt{+ Id : int}\\
    \texttt{+ DoctorId : int}\\
    \texttt{+ WeekStartDate : DateOnly}\\
    \texttt{+ PatientName : string}\\
    \texttt{+ Email : string}\\
    \texttt{+ NotifiedAt : DateTime?}
  };

  \draw[-Stealth, thick] (u.east) -- (ts.west)
    node[midway, above, font=\tiny] {1 $\to$ 0..*};
  \draw[-Stealth, thick] (ts.south) -- (ap.north)
    node[midway, right, font=\tiny] {1 $\to$ 0..1};
  \draw[-Stealth, thick] (u.south) -- (wl.north)
    node[midway, left, font=\tiny] {1 $\to$ 0..*};

\end{tikzpicture}
\caption{Diagramme de classes -- MedicalAppointments}
\end{figure}

\section{Diagrammes de Séquences}

\subsection{Séquence -- Connexion JWT (ms-auth)}

\begin{figure}[H]
\centering
\begin{tikzpicture}[every node/.style={font=\scriptsize},
  lifeline/.style={draw, thick, dashed, medsysgray},
  msg/.style={-Stealth, thick, medsysblue},
  ret/.style={-Stealth, thick, dashed, medsysgray}]

  % Participants
  \node[draw, fill=medsyslight, minimum width=2cm, align=center] (cli) at (0,0)    {Client\\(medsys-web)};
  \node[draw, fill=entitybg, minimum width=2cm, align=center]    (ac)  at (4.5,0)  {AuthController};
  \node[draw, fill=entitybg, minimum width=2cm, align=center]    (as)  at (9,0)    {AuthService};
  \node[draw, fill=attrbg, minimum width=2cm, align=center]      (jwt) at (13,0)   {JwtService};

  % Lignes de vie
  \draw[lifeline] (0,-0.3)  -- (0,-9);
  \draw[lifeline] (4.5,-0.3) -- (4.5,-9);
  \draw[lifeline] (9,-0.3)  -- (9,-9);
  \draw[lifeline] (13,-0.3) -- (13,-9);

  % Messages
  \draw[msg] (0,-1) -- (4.5,-1) node[midway, above, font=\tiny] {POST /api/v1/auth/login \{email, password\}};
  \draw[msg] (4.5,-2) -- (9,-2) node[midway, above, font=\tiny] {login(LoginRequest)};
  \draw[msg] (9,-3) -- (9,-3.5) node[right, font=\tiny] {loadUserByUsername(email)};
  \draw[msg] (9,-3.5) -- (9,-4) node[right, font=\tiny] {checkBruteForce(email)};
  \draw[msg] (9,-4.5) -- (13,-4.5) node[midway, above, font=\tiny] {generateToken(user)};
  \draw[ret] (13,-5.2) -- (9,-5.2) node[midway, above, font=\tiny] {accessToken};
  \draw[msg] (9,-5.8) -- (13,-5.8) node[midway, above, font=\tiny] {generateRefreshToken(user)};
  \draw[ret] (13,-6.5) -- (9,-6.5) node[midway, above, font=\tiny] {refreshToken};
  \draw[ret] (9,-7.2) -- (4.5,-7.2) node[midway, above, font=\tiny] {AuthResponse \{token, refreshToken, user\}};
  \draw[ret] (4.5,-8) -- (0,-8) node[midway, above, font=\tiny] {HTTP 200 -- AuthResponse};

\end{tikzpicture}
\caption{Diagramme de séquence -- Connexion JWT}
\end{figure}

\subsection{Séquence -- Prise de rendez-vous (MedicalAppointments)}

\begin{figure}[H]
\centering
\begin{tikzpicture}[every node/.style={font=\scriptsize},
  lifeline/.style={draw, thick, dashed, medsysgray},
  msg/.style={-Stealth, thick, medsysblue},
  ret/.style={-Stealth, thick, dashed, medsysgray}]

  \node[draw, fill=medsyslight, minimum width=1.8cm, align=center] (cli) at (0,0)    {Patient\\(web)};
  \node[draw, fill=entitybg, minimum width=1.8cm, align=center]    (ac)  at (3.5,0)  {Appointments\\Controller};
  \node[draw, fill=entitybg, minimum width=1.8cm, align=center]    (as)  at (7,0)    {Appointment\\Service};
  \node[draw, fill=attrbg, minimum width=1.8cm, align=center]      (db)  at (10.5,0) {AppDbContext\\(SQL Server)};
  \node[draw, fill=medsysgreen!20, minimum width=1.8cm, align=center] (nt) at (14,0) {Notification\\Service};

  \draw[lifeline] (0,-0.3)   -- (0,-11);
  \draw[lifeline] (3.5,-0.3) -- (3.5,-11);
  \draw[lifeline] (7,-0.3)   -- (7,-11);
  \draw[lifeline] (10.5,-0.3)-- (10.5,-11);
  \draw[lifeline] (14,-0.3)  -- (14,-11);

  \draw[msg] (0,-1)    -- (3.5,-1)    node[midway, above, font=\tiny] {POST /api/appointments};
  \draw[msg] (3.5,-2)  -- (7,-2)      node[midway, above, font=\tiny] {BookAsync(dto, userId?)};
  \draw[msg] (7,-3)    -- (10.5,-3)   node[midway, above, font=\tiny] {GetSlot(timeSlotId) -- vérification status=Available};
  \draw[ret] (10.5,-3.7) -- (7,-3.7) node[midway, above, font=\tiny] {TimeSlot};
  \draw[msg] (7,-4.5)  -- (10.5,-4.5) node[midway, above, font=\tiny] {CreatePatient() si anonyme};
  \draw[msg] (7,-5.3)  -- (10.5,-5.3) node[midway, above, font=\tiny] {SaveAppointment()};
  \draw[msg] (7,-6.1)  -- (10.5,-6.1) node[midway, above, font=\tiny] {UpdateSlot(status=Reserved)};
  \draw[ret] (10.5,-6.9) -- (7,-6.9) node[midway, above, font=\tiny] {Appointment créé};
  \draw[msg] (7,-7.7)  -- (14,-7.7)   node[midway, above, font=\tiny] {SendConfirmationEmail(patient)};
  \draw[msg] (14,-8.5) -- (14,-9)     node[right, font=\tiny] {Email SMTP envoyé};
  \draw[ret] (7,-9.5)  -- (3.5,-9.5)  node[midway, above, font=\tiny] {AppointmentResponseDto};
  \draw[ret] (3.5,-10.3) -- (0,-10.3) node[midway, above, font=\tiny] {HTTP 201 Created};

\end{tikzpicture}
\caption{Diagramme de séquence -- Prise de rendez-vous}
\end{figure}

\subsection{Séquence -- Consultation du dossier médical}

\begin{figure}[H]
\centering
\begin{tikzpicture}[every node/.style={font=\scriptsize},
  lifeline/.style={draw, thick, dashed, medsysgray},
  msg/.style={-Stealth, thick, medsysblue},
  ret/.style={-Stealth, thick, dashed, medsysgray}]

  \node[draw, fill=medsyslight, minimum width=1.8cm, align=center] (cli) at (0,0)   {Patient};
  \node[draw, fill=entitybg, minimum width=1.8cm, align=center]    (gw)  at (4,0)   {JwtAuthFilter};
  \node[draw, fill=entitybg, minimum width=1.8cm, align=center]    (pc)  at (8,0)   {PatientPortal\\Controller};
  \node[draw, fill=attrbg, minimum width=1.8cm, align=center]      (ps)  at (12,0)  {Patient\\Service};

  \draw[lifeline] (0,-0.3)  -- (0,-9);
  \draw[lifeline] (4,-0.3)  -- (4,-9);
  \draw[lifeline] (8,-0.3)  -- (8,-9);
  \draw[lifeline] (12,-0.3) -- (12,-9);

  \draw[msg] (0,-1) -- (4,-1) node[midway, above, font=\tiny] {GET /api/v1/patient/me/dossier (Authorization: Bearer JWT)};
  \draw[msg] (4,-2) -- (4,-2.5) node[right, font=\tiny] {validateToken(JWT)};
  \draw[msg] (4,-3) -- (4,-3.5) node[right, font=\tiny] {extractUserId(), injectHeaders};
  \draw[msg] (4,-4) -- (8,-4) node[midway, above, font=\tiny] {GET /dossier (X-User-Id: patientId)};
  \draw[msg] (8,-5) -- (12,-5) node[midway, above, font=\tiny] {getMyDossier(patientId)};
  \draw[msg] (12,-6) -- (12,-6.5) node[right, font=\tiny] {findByPatientId(), eager load all};
  \draw[ret] (12,-7) -- (8,-7) node[midway, above, font=\tiny] {DossierMedicalDTO (consultations, ordonnances, analyses...)};
  \draw[ret] (8,-7.8) -- (4,-7.8) node[midway, above, font=\tiny] {HTTP 200};
  \draw[ret] (4,-8.5) -- (0,-8.5) node[midway, above, font=\tiny] {DossierMedicalDTO JSON};

\end{tikzpicture}
\caption{Diagramme de séquence -- Consultation du dossier médical}
\end{figure}

\section{Diagramme de Composants}

\begin{figure}[H]
\centering
\begin{tikzpicture}[every node/.style={font=\scriptsize},
  comp/.style={draw, thick, rounded corners=4pt, minimum width=3.8cm,
               minimum height=1.4cm, align=center},
  db/.style={draw, thick, cylinder, shape border rotate=90,
             fill=gray!15, minimum width=2.2cm, minimum height=0.9cm,
             align=center, font=\tiny\bfseries},
  iface/.style={draw, circle, fill=medsysblue, minimum size=4pt,
                inner sep=0pt}]

  % Couche Client
  \node[comp, fill=medsysgreen!15] (web) at (6,9)
    {\textbf{medsys-web}\\React 18 + Vite\\Port 5173};
  \node[comp, fill=medsysgreen!10] (mob) at (11,9)
    {\textbf{medsys-mobile}\\React Native / Expo};

  % Services Java
  \node[comp, fill=entitybg] (auth) at (0,5.5)
    {\textbf{ms-auth}\\Spring Boot 3\\Port 8082};
  \node[comp, fill=entitybg] (pat) at (5,5.5)
    {\textbf{ms-patient-personnel}\\Spring Boot 3\\Port 8081};

  % Services .NET
  \node[comp, fill=relbg!60] (rdv) at (10,5.5)
    {\textbf{MedicalAppointments}\\ASP.NET Core 8\\Port 5000};
  \node[comp, fill=medsysgreen!20] (per) at (15,5.5)
    {\textbf{ms-personnel}\\ASP.NET Core 8\\Port var.};

  % Bases de données
  \node[db] (dbauth) at (0,3) {MySQL\\ms\_auth\_db};
  \node[db] (dbpat)  at (5,3) {MySQL\\ms\_patient\_db};
  \node[db] (dbrdv)  at (10,3) {SQL Server\\medical\_appts};
  \node[db] (dbper)  at (15,3) {MongoDB\\personnel\_db};

  % RabbitMQ
  \node[comp, fill=orange!20, minimum width=3cm, minimum height=1cm] (rmq)
    at (2.5,1) {\textbf{RabbitMQ}\\Port 5672};

  % Connexions client -> services
  \draw[-Stealth, thick, medsysblue] (web) -- (auth)
    node[midway, left, font=\tiny] {/api};
  \draw[-Stealth, thick, medsysblue] (web) -- (pat)
    node[midway, right, font=\tiny] {/api/v1};
  \draw[-Stealth, thick, orange!70!red] (web) -- (rdv)
    node[midway, right, font=\tiny] {/api/appts};
  \draw[-Stealth, thick, medsysblue] (mob) -- (auth);
  \draw[-Stealth, thick, medsysblue!50] (web) to[bend left=10] (per);

  % Services -> BDD
  \draw[-Stealth, thick, gray] (auth) -- (dbauth);
  \draw[-Stealth, thick, gray] (pat)  -- (dbpat);
  \draw[-Stealth, thick, gray] (rdv)  -- (dbrdv);
  \draw[-Stealth, thick, gray] (per)  -- (dbper);

  % RabbitMQ
  \draw[-Stealth, thick, dashed, orange!70] (auth) -- (rmq)
    node[midway, left, font=\tiny] {pub};
  \draw[-Stealth, thick, dashed, orange!70] (pat)  -- (rmq)
    node[midway, right, font=\tiny] {pub/sub};

  % Inter-service
  \draw[Stealth-Stealth, thick, dashed, medsysgray!60] (auth) -- (pat)
    node[midway, above, font=\tiny] {JWT};

\end{tikzpicture}
\caption{Diagramme de composants -- Architecture MedSys}
\end{figure}

\chapter{Description des Modules}

\section{ms-auth -- Service d'authentification}

\subsection{Responsabilités}
\begin{itemize}[itemsep=3pt]
  \item Gestion du cycle de vie des comptes utilisateurs (inscription, connexion, désactivation).
  \item Émission et validation des tokens JWT (accès 15 min + refresh 7 jours).
  \item Protection contre les attaques par force brute (\texttt{LoginRateLimiter}).
  \item Vérification email par lien tokenisé.
  \item Traçabilité de toutes les actions sensibles (audit logs).
\end{itemize}

\subsection{Endpoints principaux}

\begin{table}[H]
\centering
\renewcommand{\arraystretch}{1.3}
\begin{tabular}{|l|l|l|l|}
\hline
\rowcolor{tabhead}\color{white}\textbf{Méthode} &
\color{white}\textbf{URL} &
\color{white}\textbf{Accès} &
\color{white}\textbf{Description} \\
\hline
POST & \code{/api/v1/auth/register} & Public & Inscription patient \\
\rowcolor{roweven}
POST & \code{/api/v1/auth/login} & Public & Connexion + JWT \\
POST & \code{/api/v1/auth/refresh} & Public & Rafraîchir le token \\
\rowcolor{roweven}
POST & \code{/api/v1/auth/forgot-password} & Public & Demande reset mdp \\
POST & \code{/api/v1/auth/reset-password} & Public & Réinitialiser mdp \\
\rowcolor{roweven}
POST & \code{/api/v1/auth/change-password} & Connecté & Changer mdp \\
GET  & \code{/api/v1/auth/me} & Connecté & Profil courant \\
\rowcolor{roweven}
POST & \code{/api/v1/admin/personnel} & ADMIN & Créer compte personnel \\
\hline
\end{tabular}
\caption{Endpoints ms-auth}
\end{table}

\section{ms-patient-personnel -- Dossier médical}

\subsection{Responsabilités}
\begin{itemize}[itemsep=3pt]
  \item Gestion complète des dossiers médicaux électroniques.
  \item Portail patient : consultation, messagerie, documents, QR Code, export PDF.
  \item Dashboard direction : statistiques, accès tous patients, export PDF individuel.
  \item Proxy vers MedicalAppointments pour l'affichage des RDV.
\end{itemize}

\subsection{Structure du dossier médical}

\begin{table}[H]
\centering
\renewcommand{\arraystretch}{1.3}
\begin{tabular}{|l|l|}
\hline
\rowcolor{tabhead}\color{white}\textbf{Section} &
\color{white}\textbf{Contenu} \\
\hline
Antécédents & Médicaux, chirurgicaux, familiaux, allergies (avec sévérité) \\
\rowcolor{roweven}
Consultations & Date, médecin, motif, diagnostic, traitement, constantes vitales \\
Ordonnances & Médicaments, dosages, durées, instructions \\
\rowcolor{roweven}
Analyses labo & Type, laboratoire, statut, résultats \\
Radiologies & Type examen, description, conclusion \\
\rowcolor{roweven}
Hospitalisations & Service, dates entrée/sortie, motif \\
Certificats médicaux & Type, contenu, date émission \\
\rowcolor{roweven}
Documents uploadés & Ordonnances scannées, radios, analyses (PDF/images) \\
\hline
\end{tabular}
\caption{Structure du dossier médical électronique}
\end{table}

\section{MedicalAppointments -- Prise de rendez-vous}

\subsection{Responsabilités}
\begin{itemize}[itemsep=3pt]
  \item Gestion des créneaux horaires par les médecins (création unitaire ou en masse, blocage).
  \item Réservation de rendez-vous pour les patients authentifiés \emph{ou anonymes}.
  \item Annulation avec token anonyme (pour les patients sans compte).
  \item Liste d'attente avec notification email automatique à la libération d'un créneau.
  \item Notifications temps réel via SignalR (\code{/appointmenthub}).
\end{itemize}

\subsection{Statuts d'un rendez-vous}

\begin{table}[H]
\centering
\renewcommand{\arraystretch}{1.3}
\begin{tabular}{|l|l|l|}
\hline
\rowcolor{tabhead}\color{white}\textbf{Statut} &
\color{white}\textbf{Description} &
\color{white}\textbf{Transition depuis} \\
\hline
Confirmed & RDV confirmé par défaut à la réservation & (initial) \\
\rowcolor{roweven}
Pending & En attente de confirmation manuelle & Confirmed \\
CancelledByPatient & Annulé par le patient & Confirmed, Pending \\
\rowcolor{roweven}
CancelledByDoctor & Annulé par le médecin & Confirmed, Pending \\
CancelledBySecretary & Annulé par la secrétaire & Confirmed, Pending \\
\rowcolor{roweven}
Completed & Consultation effectuée & Confirmed \\
NoShow & Patient absent & Confirmed \\
\hline
\end{tabular}
\caption{Cycle de vie d'un rendez-vous}
\end{table}

\subsection{Règle de pénalité patient}
Si un patient annule ou est absent plus de 3 fois, le système enregistre
une date de \texttt{PenaltyUntil} qui bloque temporairement ses réservations.
Le compteur \texttt{CancelCount} est incrémenté à chaque annulation.

\section{ms-personnel -- Gestion RH}

\subsection{Responsabilités}
\begin{itemize}[itemsep=3pt]
  \item Gestion des profils du personnel médical et paramédical.
  \item Planifications des gardes et disponibilités.
  \item Gestion des absences (congés, maladie, formation).
  \item Organisation en équipes avec responsable.
  \item Demandes de modification de planning.
\end{itemize}

\subsection{Types de personnel gérés}

\begin{itemize}
  \item Médecins (avec spécialité, créneaux de consultation)
  \item Infirmiers
  \item Aides-soignants
  \item Brancardiers
  \item Secrétaires médicales
  \item Directeurs de service
\end{itemize}

\section{medsys-web -- Interface utilisateur}

\subsection{Tableaux de bord par rôle}

\begin{table}[H]
\centering
\renewcommand{\arraystretch}{1.3}
\begin{tabular}{|l|l|}
\hline
\rowcolor{tabhead}\color{white}\textbf{Rôle} &
\color{white}\textbf{Fonctionnalités accessibles} \\
\hline
PATIENT & Accueil, Dossier médical, Prendre RDV, Mes RDV,
  Documents, Messagerie, QR/PDF, Profil \\
\rowcolor{roweven}
MÉDECIN / PERSONNEL & Liste patients, Dossier complet par patient \\
DIRECTEUR & Vue d'ensemble (KPIs), Patients, Médecins,
  Comptes, RDV \& Planning, Stats RDV \\
\rowcolor{roweven}
ADMIN & Gestion des comptes (créer, activer/désactiver) \\
\hline
\end{tabular}
\caption{Fonctionnalités par rôle dans medsys-web}
\end{table}

\subsection{Flux d'authentification frontend}

\begin{enumerate}[itemsep=3pt]
  \item Le patient se connecte via \textbf{LandingPage} (split login).
  \item Le token JWT est stocké en \code{sessionStorage} sous la clé \code{medsys\_token}.
  \item \textbf{AuthContext} décode le token et expose le rôle à tous les composants.
  \item React Router redirige vers le dashboard approprié selon le rôle.
  \item À l'expiration du token (15 min), un refresh automatique est tenté.
\end{enumerate}

\subsection{Intégration MedicalAppointments}
La section \textbf{Prendre RDV} du dashboard patient utilise directement
l'API MedicalAppointments via le proxy Vite. Le flux est :

\begin{enumerate}[itemsep=3pt]
  \item \textbf{Étape 1} -- Sélection du service médical (\code{GET /api/services}).
  \item \textbf{Étape 2} -- Sélection du médecin (\code{GET /api/services/\{id\}/doctors}).
  \item \textbf{Étape 3} -- Vue semaine avec créneaux (\code{GET /api/slots?doctorId=\&weekStart=}).
  \item \textbf{Confirmation} -- Modal de confirmation avec motif, réservation
    (\code{POST /api/appointments}), affichage du token d'annulation si anonyme.
  \item \textbf{Liste d'attente} -- Si semaine complète : \code{POST /api/waiting-list}.
\end{enumerate}

\chapter{Conclusion et Perspectives}

\section{Bilan du projet}

MedSys constitue une réponse complète aux défis de la transformation numérique
hospitalière. L'architecture microservices adoptée permet d'atteindre les objectifs
initiaux tout en offrant une base solide pour les évolutions futures.

\begin{table}[H]
\centering
\renewcommand{\arraystretch}{1.4}
\begin{tabular}{|l|c|l|}
\hline
\rowcolor{tabhead}\color{white}\textbf{Objectif} &
\color{white}\textbf{Statut} &
\color{white}\textbf{Module concerné} \\
\hline
Authentification sécurisée JWT + RBAC & \textcolor{medsysgreen}{\bfseries Réalisé} & ms-auth \\
\rowcolor{roweven}
Dossier médical électronique complet & \textcolor{medsysgreen}{\bfseries Réalisé} & ms-patient-personnel \\
Prise de RDV en ligne (anonyme + connecté) & \textcolor{medsysgreen}{\bfseries Réalisé} & MedicalAppointments \\
\rowcolor{roweven}
Liste d'attente automatique & \textcolor{medsysgreen}{\bfseries Réalisé} & MedicalAppointments \\
Notifications email + temps réel & \textcolor{medsysgreen}{\bfseries Réalisé} & MedicalAppointments \\
\rowcolor{roweven}
Gestion du personnel hospitalier & \textcolor{medsysgreen}{\bfseries Réalisé} & ms-personnel \\
Tableau de bord direction & \textcolor{medsysgreen}{\bfseries Réalisé} & medsys-web \\
\rowcolor{roweven}
Export PDF dossier médical & \textcolor{medsysgreen}{\bfseries Réalisé} & ms-patient-personnel \\
QR Code dossier patient & \textcolor{medsysgreen}{\bfseries Réalisé} & ms-patient-personnel \\
\rowcolor{roweven}
Statistiques RDV direction & \textcolor{medsysgreen}{\bfseries Réalisé} & medsys-web \\
\hline
\end{tabular}
\caption{Récapitulatif des objectifs atteints}
\end{table}

\section{Points forts de l'architecture}

\begin{itemize}[itemsep=6pt]
  \item \textbf{Hétérogénéité technologique maîtrisée} : le projet démontre qu'il
    est possible de faire coexister des technologies Java (Spring Boot) et .NET
    (ASP.NET Core) au sein d'un même système cohérent, chaque service utilisant
    la stack la mieux adaptée à son domaine fonctionnel.

  \item \textbf{Isolation des données} : le principe \textit{database per service}
    garantit l'autonomie de chaque microservice — une migration de schéma ou une
    panne base de données n'affecte qu'un seul service.

  \item \textbf{Scalabilité horizontale} : chaque service peut être répliqué
    indépendamment selon la charge (ex. : multiplier les instances de
    MedicalAppointments lors des pics de prises de RDV).

  \item \textbf{Découplage via RabbitMQ} : les événements asynchrones
    (\code{patient.registered}, \code{appointment.created}) permettent aux services
    de communiquer sans dépendance temporelle, renforçant la résilience.

  \item \textbf{Temps réel avec SignalR} : la notification instantanée des créneaux
    libérés améliore significativement l'expérience utilisateur en évitant les
    rechargements manuels.
\end{itemize}

\section{Limites et points d'amélioration}

\begin{infobox}[Points d'amélioration identifiés]
\begin{itemize}[itemsep=4pt]
  \item \textbf{API Gateway} : en l'état, le frontend Vite proxy joue le rôle
    de routeur. En production, un vrai API Gateway (Kong, Traefik, NGINX) serait
    nécessaire pour la gestion centralisée des auth, rate-limiting, SSL/TLS et
    load balancing.

  \item \textbf{Service Discovery} : l'adressage des services est actuellement
    statique (ports codés en dur). Un registre de services (Consul, Eureka) ou
    un orchestrateur (Kubernetes) permettrait l'auto-découverte et la
    haute disponibilité.

  \item \textbf{Observabilité} : intégrer une stack de monitoring complète
    (Prometheus + Grafana pour les métriques, ELK/Loki pour les logs centralisés,
    Jaeger pour le tracing distribué).

  \item \textbf{Tests automatisés} : augmenter la couverture de tests unitaires
    (JUnit, xUnit) et d'intégration (Testcontainers), et mettre en place un
    pipeline CI/CD (GitHub Actions, GitLab CI).

  \item \textbf{Application mobile} : le frontend React Native (Expo) mentionné
    dans le stack technologique reste à connecter aux microservices.

  \item \textbf{Internationalisation} : préparer le système pour une utilisation
    multilingue (fr/ar/en) en vue d'un déploiement dans des établissements à
    clientèle internationale.
\end{itemize}
\end{infobox}

\section{Perspectives d'évolution}

\subsection{Court terme}

\begin{itemize}[itemsep=4pt]
  \item Déploiement Dockerisé complet avec \code{docker-compose.yml} unifié
    orchestrant tous les services, bases de données et RabbitMQ.
  \item Mise en place d'un reverse proxy NGINX en front de tous les services
    pour remplacer le proxy Vite en production.
  \item Intégration d'un système de paiement en ligne pour les consultations
    privées (Stripe, PayPal).
\end{itemize}

\subsection{Moyen terme}

\begin{itemize}[itemsep=4pt]
  \item Migration vers Kubernetes pour l'orchestration des conteneurs et la
    gestion automatique de la scalabilité.
  \item Implémentation d'un service de téléconsultation via WebRTC intégré
    dans le tableau de bord médecin.
  \item Développement du module de facturation et de gestion des assurances.
  \item Application mobile React Native complète pour les patients (iOS / Android).
\end{itemize}

\subsection{Long terme}

\begin{itemize}[itemsep=4pt]
  \item \textbf{Intelligence Artificielle} : aide à la décision médicale,
    détection d'interactions médicamenteuses, prédiction des pics d'affluence.
  \item \textbf{Interopérabilité HL7/FHIR} : conformité aux standards
    internationaux d'échange de données médicales pour l'intégration avec
    d'autres systèmes hospitaliers.
  \item \textbf{Blockchain} : traçabilité inaltérable des accès au dossier
    médical pour une conformité RGPD renforcée.
\end{itemize}

\section{Conclusion générale}

\begin{tcolorbox}[enhanced, colback=medsyslight, colframe=medsysdark,
  fonttitle=\bfseries, title={MedSys -- Une architecture prête pour la production},
  leftrule=5pt, arc=4pt]
MedSys démontre la viabilité d'une architecture microservices hétérogène pour
un système d'information hospitalier moderne. En combinant la robustesse de
Spring Boot, la performance d'ASP.NET Core et la réactivité de React, le
système répond aux exigences fonctionnelles tout en maintenant une séparation
claire des responsabilités.

\vspace{0.4cm}
La modélisation Merise garantit une structure de données rigoureuse et normalisée,
tandis que les diagrammes UML offrent une vision comportementale et structurelle
complète du système. Ce double niveau de modélisation facilite la maintenance,
l'onboarding de nouveaux développeurs et les évolutions futures.

\vspace{0.4cm}
Avec une sécurité transversale (JWT, RBAC, audit logs, brute-force protection),
des communications asynchrones via RabbitMQ et des notifications temps réel
via SignalR, MedSys pose les fondations d'une plateforme hospitalière numérique
robuste, scalable et prête pour un déploiement en environnement de production.
\end{tcolorbox}



\end{document}
